% Source: https://artofproblemsolving.com/community/c3h1621106_eh_some_more_geometry
% Stable identity: aops-topic-1621106-q01
The side of the square is $\sqrt{108}=6\sqrt3$. On each side, the circle cuts off a chord of length $6\sqrt3-2\left(3\sqrt3-3\right)=6$, centred at the midpoint $F$ of the side. Let the chord on the right side be $XY$.

\begin{asy}
usepackage("amsmath");
size(6cm);
defaultpen(fontsize(10pt));
real h=3*sqrt(3);pair O=(0,0);
for(int k=0;k<4;++k){ fill(rotate(90*k)*((h,3)--(h,h)--(3,h)--arc(O,6,60,30)--cycle),lightgray); }
draw((-h,-h)--(h,-h)--(h,h)--(-h,h)--cycle,linewidth(0.9));draw(circle(O,6),linewidth(0.9));
pair F=(h,0),X=(h,3),Y=(h,-3);
draw(O--X,blue+dashed);draw(O--Y,blue+dashed);draw(O--F,blue+dashed);
draw(F+(-0.5,0)--F+(-0.5,0.5)--F+(0,0.5),blue);
dot(O);label("$O$",O,W);dot(X);dot(Y);label("$X$",X,NE);label("$Y$",Y,SE);label("$F$",F,E);
label("$6$",(X+O)/2,NW,blue);label("$3\sqrt3$",(2.0,0),N,blue);label("$3$",(h,1.5),W,blue);
label("\mbox{}",(min(currentpicture,user=true)+max(currentpicture,user=true))/2,black); // keeps the text after the figure black when LaTeX typesets the labels inline
\end{asy}

\textbf{Radius.} $OF=\dfrac{6\sqrt3}{2}=3\sqrt3$ and $FX=3$, so $OX=\sqrt{27+9}=6$.

\textbf{Area of the circle inside the square.} In the right triangle $OFX$, $FX=\dfrac{1}{2}OX$, so $\angle FOX=30^\circ$ and $\angle XOY=60^\circ$. The part of the circle beyond side $XY$ is a circular segment with area
\[
\dfrac{60}{360}\pi\cdot6^2-\dfrac{\sqrt3}{4}\cdot6^2=6\pi-9\sqrt3 ,
\]
(sector minus the equilateral triangle $OXY$ of side $6$). The four segments lie outside the square, so the part of the circle inside the square has area
\[
36\pi-4\left(6\pi-9\sqrt3\right)=12\pi+36\sqrt3 .
\]
\textbf{Shaded area.} The shaded regions are the rest of the square:
\[
108-\left(12\pi+36\sqrt3\right)=\boxed{108-36\sqrt3-12\pi}\approx7.95\text{ cm}^2 .
\]
